% Бүлэг 4: Хэрэгжүүлэлтийн төлөвлөгөө ба арга зүйн шинжилгээ

\chapter{Хөгжил ба хэрэгжүүлэлтийн төлөвлөгөө}
\label{Chapter4}

%-------------------------------------------------------------------------------
\section{Системийн хэрэгжүүлэлтийн ерөнхий стратеги}

Энэхүү систем нь өндөр функциональ шаардлагатай олон давхаргат архитектур (Multi-layer Architecture) болгон хөгжүүлэгдэхээр төлөвлөгдсөн. Хэрэгжүүлэлтийн явцад аль болох түргэн үр дүнд хүрэх (Agile Development), чанарын баталгаа (Quality Assurance), болон уялдаа холбоотой үйл ажиллагааг эхлүүлэх (Continuous Integration) зарчмуудыг баримтална.

%-------------------------------------------------------------------------------
\section{Технологийн сонголт ба үндэслэл}

\subsection{Backend технологи}

Системийн үндсэн сервер талын хэрэгжүүлэлтэд дараах технологиуд ашигладаг:

\begin{itemize}
    \item \textbf{Python (Flask/FastAPI):} AI загвар болон аудио сигналын боловсруулалтыг Python-д хэрэгжүүлэх нь стандарт байдаг. TensorFlow, PyTorch зэрэг машин сурах сангууд Python-д үндэслэгдсэн бүтэцтэй байдаг. FastAPI нь асинхрон хүсэлтийг (async requests) сайн дэмждэг тул нэгэн зэрэг олон хэрэглэгчтэй үйлчлэхэд тохиромжтой.
    
    \item \textbf{PostgreSQL:} Реляционал өгөгдлийн сан болгон сонгосон. Өгөгдлийн нөхцөл бүрхэлт (ACID compliance), индексжүүлэлт, болон JSONB өгөгдлийн төрлийг дэмждэг учраас вектор өгөгдлийн сангийн шаардлагатай нийцэж байна.
    
    \item \textbf{AWS S3:} Аудио файл, загварын жинс, болон сургалтын өгөгдлийн сан (vector database)-ийг хадгалахдаа AWS S3-ыг ашигладаг. S3 нь өндөр найдвартай байдал (99.99\% availability), масштабшилт (scalability), болон доогуур төлбөр (cost-effective pricing) үзүүлэлтээр сайн үнэлэгдэж байна.
\end{itemize}

\subsection{Frontend технологи}

Хэрэглэгчийн интерфейсийн хэрэгжүүлэлтэд:

\begin{itemize}
    \item \textbf{Next.js:} React фреймворкт дүрмэлэгдэлтэй (server-side rendering), статик үүсгэлт (static generation), болон API route-уудыг дэмждэг. Next.js-ийн SEO дэмжлэг, өндөр гүйцэтгэл, болон хялбар deployment нь бидний хэрэгцээтэй нийцэж байна.
    
    \item \textbf{TypeScript:} Төрлийн аюулгүй байдал (type safety) болон код нотолгоо (code documentation) үзүүлэлтээр давуу. Том төслүүдэд заль сохлоо бууруулдаг.
    
    \item \textbf{TailwindCSS:} Хурдан UI хөгжүүлэлт, нийцтэй загвар (consistent design), болон draft нь нүүлгэн шилжүүлэх (responsive design) үзүүлэлтээр дотроо байдаг.
    
    \item \textbf{Web Audio API:} Браузер дотор аудио боловсруулалт (real-time waveform visualization, playback control) хийхэд ашигладаг.
\end{itemize}

\subsection{AI/ML технологи}

\begin{itemize}
    \item \textbf{TensorFlow/Keras:} CNN загвар сургалт болон инференс.
    
    \item \textbf{Google Magenta:} MusicVAE загварын сургалт. Эсхүл pytorch-d суурилсан Music Transformer.
    
    \item \textbf{LangChain + OpenAI/Ollama:} RAG системийн шойлго хүсэлтийн хэрэгжүүлэлт. Сээр (locally) туршилт хийхдээ Ollama (open-source LLM)-ийг, үйлдвэрт үгэй жалган мэдээллийг сайлахдаа OpenAI API эсхүл зөөнхөн туршилтын өгөгдөл дээр хязгаарлан хэрэглэж болно.
    
    \item \textbf{Librosa + Scipy:} Аудио сигналын боловсруулалт, спектрограмм үүсгэлт.
    
    \item \textbf{Chroma / FAISS:} Вектор өгөгдлийн сан RAG системийн нэмэлт мэдлэгийг хадгалахдаа.
\end{itemize}

%-------------------------------------------------------------------------------
\section{Системийн архитектурын нарийвчилсан загвар}

\subsection{Микросервис архитектур}

Системийг масштабшилтыг сайтар хийхийн тулд микросервис архитектур дагуулан хөгжүүлнэ:

\begin{enumerate}
    \item \textbf{API Gateway:} Бүх хүсэлтийг дүр сэтгэл, rate limiting, болон баталгаажуулалтын (authentication) үүргүүлэлт.
    
    \item \textbf{Auth Service:} Хэрэглэгчийн бүртгэл, нэвтрэлт, токен удирдалт.
    
    \item \textbf{Audio Processing Service:} Аудио файлын сургалт, спектрограмм үүсгэлт, CNN загварын инференс.
    
    \item \textbf{Melody Generation Service:} MusicVAE загварын инференс, MIDI үүсгэлт.
    
    \item \textbf{Chatbot Service:} RAG системийн хэрэгжүүлэлт, вектор өгөгдлийн сантай холбоо.
    
    \item \textbf{Marketplace Service:} VST плагин, Sample Pack-ийн борлуулалт, сэтгэгдэл удирдалт.
    
    \item \textbf{Database Service:} PostgreSQL с хэвлэл.
\end{enumerate}

%-------------------------------------------------------------------------------
\section{Хэрэгжүүлэлтийн үе шатны төлөвлөгөө}

Төслийг дараах 6 үе шатаар хэрэгжүүлнэ. Нийт хугацаа: 25-35 ажлын өдөр.

\subsection{Үе шат 0: Төслийн суурь ба инфраструктур (1-2 өдөр)}

\textbf{Зорилго:} Хөгжүүлэлтийн орчин, хувилбарын удирдалт, ба суут сайрүүлэлтийн хэрэгслүүдийг бэлдэх.

\begin{itemize}
    \item Git репозитор үүсгэлт (\texttt{GitHub/GitLab})
    \item \texttt{.gitignore}, \texttt{.env} файл заташ
    \item Python virtual environment (\texttt{venv}) ба Node.js ажлын орчин суулгалт
    \item Төслийн каталог бүтэц үүсгэлт:
    \begin{lstlisting}[language=bash]
    Diploma/
    ├── Backend/
    │   ├── app/
    │   │   ├── __init__.py
    │   │   ├── routes/
    │   │   ├── models/
    │   │   ├── services/
    │   │   └── utils/
    │   ├── ml_models/
    │   │   ├── cnn_mixer.py
    │   │   ├── musicvae.py
    │   │   └── rag_chatbot.py
    │   ├── config/
    │   ├── requirements.txt
    │   └── main.py
    ├── Frontend/
    │   ├── src/
    │   │   ├── components/
    │   │   ├── pages/
    │   │   ├── services/
    │   │   ├── hooks/
    │   │   ├── types/
    │   │   └── styles/
    │   ├── public/
    │   ├── next.config.js
    │   ├── tsconfig.json
    │   └── package.json
    ├── Database/
    │   ├── schema.sql
    │   ├── migrations/
    │   └── seeds/
    └── README.md
    \end{lstlisting}
    
    \item Docker-ийн файл (\texttt{Dockerfile}) болон композит (\texttt{docker-compose.yml}) суулгалт
    \item CI/CD пайпланы анхан бүрхүүлэлт (GitHub Actions)
\end{itemize}

\subsection{Үе шат 1: Backend суурь болон өгөгдлийн сан (3-5 өдөр)}

\textbf{Зорилго:} FastAPI сервер ба PostgreSQL өгөгдлийн сан суулгалт, API endpoint-уудын бүтэц үүсгэлт.

\subsubsection{Үе шатын даалгавар:}

\begin{enumerate}
    \item \textbf{PostgreSQL өгөгдлийн сан үүсгэлт:}
    \begin{itemize}
        \item Хэрэглэгчийн хүснэгт (\texttt{users}): id, email, password\_hash, full\_name, profile\_data
        \item Аудио файлын хүснэгт (\texttt{audio\_files}): id, user\_id, filename, duration, file\_path, upload\_date
        \item Миксинг анализын хүснэгт (\texttt{mixing\_analysis}): id, audio\_id, eq\_recommendations, compression\_settings, quality\_score, analyzed\_date
        \item Мелодийн өөрчлөлт (\texttt{melody\_variations}): id, midi\_file\_id, variation\_index, midi\_data, created\_date
        \item Мэдлэгийн суурь (\texttt{knowledge\_base}): id, category, content, embedding\_vector, source, verified\_by
        \item Дэлгүүрийн зүйл (\texttt{marketplace\_items}): id, seller\_id, title, description, price, category, file\_path
        \item Сэтгэгдэлүүд (\texttt{user\_feedback}): id, user\_id, content, rating, created\_date
    \end{itemize}
    
    \item \textbf{FastAPI ба CORS суулгалт:}
    \begin{lstlisting}[language=Python]
    pip install fastapi uvicorn sqlalchemy psycopg2-binary python-dotenv pydantic
    \end{lstlisting}
    
    \item \textbf{Core API endpoint-ууд:}
    \begin{itemize}
        \item \texttt{POST /auth/register} - Хэрэглэгчийн бүртгэл
        \item \texttt{POST /auth/login} - Нэвтрэлт (JWT token)
        \item \texttt{GET /auth/profile} - Хэрэглэгчийн профайл авах
        \item \texttt{POST /audio/upload} - Аудио файл байршуулах (multipart/form-data)
        \item \texttt{GET /audio/\{id\}} - Аудио файлын мета мэдээлэл авах
        \item \texttt{DELETE /audio/\{id\}} - Аудио файл устгах
        \item \texttt{GET /marketplace} - Дэлгүүрийн зүйлсийн жагсаалт авах (pagination)
        \item \texttt{POST /marketplace/create} - Дэлгүүрт шинэ зүйл нэмэх
    \end{itemize}
    
    \item \textbf{ORM загвар ба схем үүсгэлт} (SQLAlchemy):
    \begin{lstlisting}[language=Python]
    from sqlalchemy import Column, Integer, String, Text, DateTime
    from sqlalchemy.ext.declarative import declarative_base
    from datetime import datetime
    
    Base = declarative_base()
    
    class User(Base):
        __tablename__ = "users"
        id = Column(Integer, primary_key=True)
        email = Column(String, unique=True, index=True)
        password_hash = Column(String)
        full_name = Column(String)
        created_at = Column(DateTime, default=datetime.utcnow)
    \end{lstlisting}
    
    \item \textbf{Баталгаажуулалт ба нэвтрэх (Authentication):}
    \begin{itemize}
        \item bcrypt-ийн үед нууц үгийн хеш хийх
        \item JWT token-ийн үүсгэлт ба шалгалт
        \item Refresh token механизм
    \end{itemize}
    
    \item \textbf{Файлын сано ба AWS S3 холбоо:}
    \begin{lstlisting}[language=Python]
    import boto3
    s3_client = boto3.client('s3', region_name='us-east-1')
    s3_client.upload_file('local_file', 'bucket_name', 'object_key')
    \end{lstlisting}
\end{enumerate}

\subsection{Үе шат 2: AI загвар хаяглалт (5-7 өдөр)}

\textbf{Зорилго:} CNN, MusicVAE, RAG загварыг сервертэй холбоно.

\subsubsection{Үе шатын даалгавар:}

\begin{enumerate}
    \item \textbf{CNN Audio Mixer загвар:}
    \begin{itemize}
        \item Librosa ашиглан аудио спектрограмм үүсгэлт:
        \begin{lstlisting}[language=Python]
        import librosa
        import numpy as np
        
        y, sr = librosa.load(audio_path)
        S = librosa.feature.melspectrogram(y=y, sr=sr)
        S_db = librosa.power_to_db(S, ref=np.max)
        \end{lstlisting}
        
        \item Урьдчилан сургалтан CNN загварыг ачаалж эсхүл сургалт хийх
        \item Спектрограммын нарвачилсан анализ (feature extraction)
        \item EQ ба компрессорын тохиргооны саналыг буцаах
        \item REST API endpoint: \texttt{POST /analyze/mixing}
    \end{itemize}
    
    \item \textbf{MusicVAE мелодийн үүсгэлт:}
    \begin{itemize}
        \item Google Magenta MusicVAE загварыг ачаалж сургалт хийх
        \item MIDI файлын оролт авах ба спектроскопит дүрслэлд хөрвүүлэх
        \item VAE латент орон зайд кодлох (encoding)
        \item 3-5 өөрчлөлт үүсгэх (interpolation, variation)
        \item MIDI файл үүсгэж буцаах
        \item REST API endpoint: \texttt{POST /generate/melody}
    \end{itemize}
    
    \item \textbf{RAG Chatbot систем:}
    \begin{itemize}
        \item LangChain ба ChromaDB суулгалт:
        \begin{lstlisting}[language=Python]
        pip install langchain chromadb openai python-dotenv
        \end{lstlisting}
        
        \item Мэдлэгийн суурь (knowledge base)-ийг embedding болгон хувиргалт
        \item Vector database-д сүүлчүүлэх
        \item Хэрэглэгчийн асуултад хайлт ба RAG санала үүсгэлт
        \item Hallucination угаасны механизм (confidence threshold)
        \item REST API endpoint: \texttt{POST /chat/ask}
        \item Вектор өгөгдлийн сан анхлан сургалт
    \end{itemize}
\end{enumerate}

\subsection{Үе шат 3: Frontend хөгжүүлэлт (4-6 өдөр)}

\textbf{Зорилго:} Next.js-ийг ашиглан сүүтэй интерфэйсийг хөгжүүлэх.

\subsubsection{Үе шатын даалгавар:}

\begin{enumerate}
    \item \textbf{Next.js төслийг үүсгэлт:}
    \begin{lstlisting}[language=bash]
    npx create-next-app@latest frontend --typescript --tailwind
    \end{lstlisting}
    
    \item \textbf{Баталгаажуулалт (Authentication) хуудас:}
    \begin{itemize}
        \item RegisterPage: нэр, имэйл, нууц үгийн маягтаа
        \item LoginPage: имэйл ба нууц үгийн баталгаажуулалт
        \item JWT token-ийг localStorage-д хадгалах
        \item Protected routes ба middleware
    \end{itemize}
    
    \item \textbf{Dashboard хуудас:}
    \begin{itemize}
        \item Хэрэглэгчийн профайлын товчоо
        \item Сүүлийн аудио анализуудын түүх
        \item Хүрэх мэдээлэл
    \end{itemize}
    
    \item \textbf{Аудио Анализ хуудас:}
    \begin{itemize}
        \item Drag-and-drop файл байршуулсан виджет
        \item Спектрограмм дүрслэлт (Plotly эсхүл Chart.js)
        \item EQ ба компрессорын сөөлөмжийг харуулах
        \item Дүн хүргүүлэх 0-100 үзүүлэлт
        \item Урьдчилан сургалтан анализын явцыг харуулах (progress bar)
    \end{itemize}
    
    \item \textbf{Мелодийн үүсгэлт хуудас:}
    \begin{itemize}
        \item MIDI файл байршуулах
        \item 3-5 өөрчлөлт сонголтуудыг харуулах
        \item MIDI дуу таталцаа (playback)
        \item MIDI файл татан авах (download)
    \end{itemize}
    
    \item \textbf{AI Chatbot интерфэйс:}
    \begin{itemize}
        \item Уголдсон чатын хэлэлцээ
        \item Категорийн сонголт (техник, сургалт, онол)
        \item Сөөлөмжийн үнэлгээ (feedback rating)
    \end{itemize}
    
    \item \textbf{Дэлгүүр (Marketplace) хуудас:}
    \begin{itemize}
        \item VST плагин болон Sample Pack-ийн жагсаалт
        \item Хайлт ба сүлүүлэлт
        \item Борлуулалтын түүх
    \end{itemize}
    
    \item \textbf{Өнгөний өөрчлөлт (Dark Mode):}
    \begin{itemize}
        \item TailwindCSS-ийн dark mode дэмжлэг
        \item Хэрэглэгчийн сонголтыг localStorage-д хадгалах
    \end{itemize}
\end{enumerate}

\subsection{Үе шат 4: Интеграци ба туршилт (3-4 өдөр)}

\textbf{Зорилго:} Frontend ба Backend-ийн холбоо, чанарын баталгаа.

\begin{enumerate}
    \item \textbf{API интеграци туршилт:}
    \begin{itemize}
        \item Postman эсхүл Insomnia ашиглан API endpoint-уудыг туршилт хийх
        \item Frontend сүлүүлэлт, лоадинг, алдаа удирдалт
        \item Backend алдааны сэтгэгдлийн сөлбөрт
    \end{itemize}
    
    \item \textbf{Unit testing:}
    \begin{itemize}
        \item Backend: pytest ашиглан функцийн туршилт
        \item Frontend: Jest ба React Testing Library
    \end{itemize}
    
    \item \textbf{Load testing:}
    \begin{itemize}
        \item Apache JMeter эсхүл Locust ашиглан нэгэн зэрэл 100-1000 хэрэглэгч симуляци
        \item CNN инференс хугацаа 10 сек-ийн дор байгаа эсэх шалгалт
    \end{itemize}
    
    \item \textbf{Аюулгүй байдлын шалгалт:}
    \begin{itemize}
        \item SQL injection эргүүлэх
        \item XSS (Cross-Site Scripting) хөндөлтийн сөлбөр
        \item CSRF (Cross-Site Request Forgery) дэмжлэг
        \item Rate limiting (нэг IP хаяг дээр минут тутамд 60 хүсэлт)
    \end{itemize}
    
    \item \textbf{Performance optimization:}
    \begin{itemize}
        \item Database query optimize хийх (indexing, query planning)
        \item Frontend bundle size багасгалт (code splitting, lazy loading)
        \item CDN суулгалт (CSS, JS файлын хагалалтыг хурдасгалт)
    \end{itemize}
\end{enumerate}

\subsection{Үе шат 5: Deployment ба документ (2-3 өдөр)}

\textbf{Зорилго:} Үйлдвэрт байршуулах болон сэргэх документ.

\begin{enumerate}
    \item \textbf{AWS дээр Deploy хийх:}
    \begin{itemize}
        \item RDS-д PostgreSQL үүсгэлт
        \item EC2 эсхүл ECS-д Backend суулгалт
        \item S3 дээр аудио файлын сано
        \item CloudFront CDN суулгалт
        \item Route53 DNS тохиргоо
    \end{itemize}
    
    \item \textbf{Frontend deployment:}
    \begin{itemize}
        \item Vercel эсхүл Netlify дээр Next.js deploy хийх
        \item Environment variables тохиргоо
        \item Custom domain суулгалт
    \end{itemize}
    
    \item \textbf{Docker контейнеризация:}
    \begin{lstlisting}[language=Dockerfile]
    FROM python:3.11-slim
    WORKDIR /app
    COPY requirements.txt .
    RUN pip install -r requirements.txt
    COPY . .
    CMD ["uvicorn", "main:app", "--host", "0.0.0.0"]
    \end{lstlisting}
    
    \item \textbf{CI/CD Pipeline (GitHub Actions):}
    \begin{itemize}
        \item Testing, building, deploying автомат хийлэх
        \item Хүлээх шалгалт (linting), type checking
    \end{itemize}
    
    \item \textbf{Документ бичих:}
    \begin{itemize}
        \item API документ (Swagger/OpenAPI)
        \item Хэрэглэгчийн гарын авлага
        \item Админ удирдаларыг
        \item Сүүлийн ажлын тайлан ба фотограф
    \end{itemize}
\end{enumerate}

%-------------------------------------------------------------------------------
\section{Өгөгдлийн сан (Database) дизайн}

\subsection{Схемийн нарийвчилсан загвар}

\begin{table}[!htbp]
\caption{Хэрэглэгч хүснэгтийн схем}
\centering
\begin{tabular}{|l|l|l|l|}
\hline
\textbf{Баганын нэр} & \textbf{Өгөгдлийн төрөл} & \textbf{Хязгаарлалт} & \textbf{Тайлбар} \\
\hline
id & SERIAL & PRIMARY KEY & Уникум таних тэмдэг \\
\hline
email & VARCHAR(255) & UNIQUE, NOT NULL & Хэрэглэгчийн имэйл хаяг \\
\hline
password\_hash & VARCHAR(255) & NOT NULL & bcrypt хешийн нууц үг \\
\hline
full\_name & VARCHAR(255) & NOT NULL & Бүтэн нэр \\
\hline
profile\_pic & TEXT & NULL & S3 дээрх зураг холбоос \\
\hline
created\_at & TIMESTAMP & DEFAULT NOW() & Үүсгэлтийн цаг \\
\hline
updated\_at & TIMESTAMP & DEFAULT NOW() & Сүүлийн өөрчлөлтийн цаг \\
\hline
\end{tabular}
\end{table}

\begin{table}[!htbp]
\caption{Аудио файл хүснэгтийн схем}
\centering
\begin{tabular}{|l|l|l|l|}
\hline
\textbf{Баганын нэр} & \textbf{Өгөгдлийн төрөл} & \textbf{Хязгаарлалт} & \textbf{Тайлбар} \\
\hline
id & SERIAL & PRIMARY KEY & Уникум таних тэмдэг \\
\hline
user\_id & INTEGER & FOREIGN KEY & Хэрэглэгчийн ID \\
\hline
filename & VARCHAR(255) & NOT NULL & Файлын нэр \\
\hline
duration & FLOAT & & Аудиогийн хугацаа (сек) \\
\hline
file\_path & TEXT & NOT NULL & S3 дээрх холбоос \\
\hline
upload\_date & TIMESTAMP & DEFAULT NOW() & Байршуулсан цаг \\
\hline
\end{tabular}
\end{table}

%-------------------------------------------------------------------------------
\section{Хэрэгжүүлэлтийн цаг хуваарь}

\begin{table}[!htbp]
\caption{Системийн хөгжүүлэлтийн үе шатын цаг хуваарь}
\centering
\begin{tabular}{|l|c|c|c|}
\hline
\textbf{Үе шат} & \textbf{Хугацаа} & \textbf{Байцаалтын эхлэлийн дүүнээр} & \textbf{Статус} \\
\hline
Үе шат 0: Суурь & 1-2 өдөр & 1-2 өдөр & ⏳ Чухал \\
\hline
Үе шат 1: Backend & 3-5 өдөр & 3-7 өдөр & Бэлтгэл \\
\hline
Үе шат 2: AI загвар & 5-7 өдөр & 8-14 өдөр & Бэлтгэл \\
\hline
Үе шат 3: Frontend & 4-6 өдөр & 12-20 өдөр & Бэлтгэл \\
\hline
Үе шат 4: Туршилт & 3-4 өдөр & 15-24 өдөр & Бэлтгэл \\
\hline
Үе шат 5: Deploy & 2-3 өдөр & 17-27 өдөр & Бэлтгэл \\
\hline
\textbf{Нийт} & \textbf{18-27 өдөр} & - & \textbf{~1 сар} \\
\hline
\end{tabular}
\end{table}

%-------------------------------------------------------------------------------
\section{Ирээдүйн сайжруулалтын боломжууд}

Энэхүү системийг үндэс болгон дараах сайжруулалт болон өргөтгөлүүдийг хийж болно:

\begin{itemize}
    \item \textbf{Mobile Application:} React Native эсхүл Flutter ашиглан гар утас дээрх худалдан авалтын програм
    
    \item \textbf{Real-time Collaboration:} WebSocket ашиглан хэрэглэгчид нэгэн зэрэл ажилла
    
    \item \textbf{Advanced Analytics:} Хэрэглэгчийн ашиглалтын дүн шинжилгээ, сурчилгаа
    
    \item \textbf{Community Features:} Хэрэглэгчийн бүтээлүүдийн хувьцаа, дүгнэлт, лайк систем
    
    \item \textbf{Music Recommendation:} Хэрэглэгчийн сонирхлын дагуу хөгжим сөглөлгөлтийн систем
\end{itemize}

%-------------------------------------------------------------------------------
\section{Эрсдлийн үнэлгээ ба шийдвэрлэх аргууд}

\begin{table}[!htbp]
\caption{Төслийн эрсдэлүүд ба хөрвүүлэлтийн стратеги}
\centering
\begin{tabular}{|p{3cm}|p{3cm}|p{2cm}|p{4cm}|}
\hline
\textbf{Эрсдэл} & \textbf{Шалтгаан} & \textbf{Магадлал} & \textbf{Хөрвүүлэлтийн арга} \\
\hline
CNN загварын Нарвачилсан байдал Бага & Сургалтын өгөгдлийн сан дутагдалтай & Өндөр & Өмнө сургалтын загвар (transfer learning) ашиглах \\
\hline
API гүйцэтгэл удаан & Олон хэрэглэгч нэгэн зэрэл & Дээд & Asynchronous processing, task queue (Celery) ашиглах \\
\hline
AWS S3 зардлын өндөр & Олон аудио файл & Дээд & Файл сөлбөр, автоматаар хуучин файлүүдийг устгалт \\
\hline
Өгөгдлийн аюулгүй байдлын эрсдэл & Сансан нэвтрэлт & Дээд & HTTPS, двалын баталгаажуулалт (2FA), rate limiting \\
\hline
\end{tabular}
\end{table}

%-------------------------------------------------------------------------------
\section{Бүлгийн дүгнэлт}

Энэ бүлгээр системийн хэрэгжүүлэлтийн ерөнхий стратеги, технологийн сонголтын үндэслэл, болон 5 үе шатаас бүрдсэн нарийвчилсан хөгжүүлэлтийн төлөвлөгөөг үзүүлэв. Next.js, FastAPI, PostgreSQL, болон AWS S3-ыг сонгосон нь сургалтын экосистемийн хэрэгцээтэй нийцэж, масштабшилт, чанарын баталгаа, болон ирээдүйн өргөтгөл хийхэд боломжтой архитектур биелүүлэхийг зорисон болно. 

Системийн хөгжүүлэлтийн явцад агил аргачлалыг баримтлан, эхлүүлэлт хийх, туршилт хийхэд үгүйцэлбэлж, итеративаар сайжруулна. Өгөгдлийн аюулгүй байдал, гүйцэтгэлийн үзүүлэлт, болон хэрэглэгчийн туршлагыг чүн төвөтэй хөгжүүлнэ.

